
\needspace{2\baselineskip}
\color{unidarkgreen}{\subsection{Комбинированный пуск по напряжению стороны НН}}
\color{black}

\begin{enumerate}[label=\arabic{section}.\arabic{subsection}.\arabic*,  align=left, labelsep=4pt, leftmargin=0pt, itemindent=57pt]

\item
Комбинированный пусковой орган по напряжению низкой и средней сторон трансформатора используется для повышения чувствительности МТЗ НН.   

\item
Цепи напряжения подключаются к ТН, установленным на шинах низшего напряжения.

\item
Внутренний комбинированный пусковой орган напряжения (КПОН), контролирует минимальное линейное напряжение и максимальное напряжение обратной последовательности низкой стороны трансформатора. Порог срабатывания пуска по минимальному линейному напряжению задается уставками \textbf{«Uл КПОН НН»}, порог срабатывания пуска по максимальному напряжению обратной последовательности задается уставками \textbf{«U2 КПОН НН»}. 

\item
Уставки КПОН приведены в таблице \ref{app_set:kpon}.

\item
Ввод в работу КПОН осуществляется независимо для каждой стороны с помощью уставок \textbf{«КПОН НН»}. Оперативный ввод/вывод КПОН производится с помощью БУ (блоков управления функциями)
\textit{«Вывод КПОН НН»}.

\item
Выбор типа пуска по напряжению соответствующей стороны осуществляется с помощью уставок \textbf{«Тип КПОН НН»}: 
\begin{itemize}
\item \textit{«Umin, U2»} - пуск по минимальному линейному напряжению или максимальному напряжению обратной последовательности;
\item \textit{«Umin»} - пуск по минимальному линейному напряжению;
\item \textit{«Внеш»} - пуск по внешнему входному сигналу «Внеш. КПОН» соответствующей стороны.
\end{itemize}

\item
Предусмотрена возможность контроля цепей напряжения каждой стороны по внутренней логике (КЦН). При появлении неисправности в цепях напряжения ТН соответствующей стороны, действие логики КПОН задается уставками \textbf{«Неиспр ТН НН»}:
\begin{itemize}
\item \textit{«Выведено»} - неисправность в цепях ТН не вызывает изменения логики КПОН соответствующей стороны;
\item \textit{«Пуск КПОН»} - при появлении неисправности в цепях ТН, КПОН соответствующей стороны всегда разрешен;
\item \textit{«Блок КПОН»} - при появлении неисправности в цепях ТН, КПОН соответствующей стороны блокируется до исчезновения неисправности ТН.
\end{itemize}

\item
Алгоритм работы КПОН НН приведен на рисунке на рисунке \ref{pic:KPON}. Алгоритмы работы КПОН НН1, КПОН НН2 выполнены аналогично.

\begin{figure}[H]
\centering
\includegraphics[width=1\textwidth,height=1\textheight,keepaspectratio]{\fbpath/06.02 ДЗАТ/081.0371 КПОН/img/KPON.pdf}
\captionof{figure}{Функциональная схема логики КПОН НН}\label{pic:KPON}
\end{figure}


\end{enumerate}